\documentclass[journal]{IEEEtran}

\usepackage{amsmath}
\usepackage{graphicx}
\usepackage{url}
\usepackage{cite}
\usepackage{array}

\hyphenation{op-tical net-works semi-conduc-tor nano-ma-chine nano-ma-chines}

\begin{document}

\title{Energy-Aware Extension of AcCoRD: Per-Nanosensor Energy Tracking in Molecular Communication Simulations}

\author{Mathis~Paccoud}

\markboth{Journal of Molecular Communication and Nanotechnology,~2026}%
{Paccoud: Energy-Aware Extension of AcCoRD}

\maketitle

\begin{abstract}
Intra-body nanosensor networks represent a promising paradigm for continuous
health monitoring and targeted drug delivery, but their practical deployment
is fundamentally constrained by energy. Existing molecular communication (MC)
simulators model diffusion, reaction kinetics, and detection of messenger
molecules, yet none integrate a per-nanosensor energy reservoir that is
depleted by detection and inter-sensor communication events. This paper
presents an energy-aware extension of AcCoRD (Actor-based Communication via
Reaction-Diffusion), a high-fidelity hybrid micro-/mesoscopic reaction-diffusion
simulator. We augment each simulated nanosensor with an individual energy
reservoir, governed by passive drain, event-based detection costs, and
inter-sensor gossip communication costs. A new biomarker identity tracking
mechanism (BmIDSet) propagates the set of detected biomarker identities
through reaction products, enabling deduplication of detection events across
the network. The extended simulator is validated on a five-region branching
arteriole scenario with 56 biomarker molecules and 100 individual nanosensors
organised in four clusters. All energy parameter values used in the simulation
are arbitrary placeholders chosen solely to exercise the model and produce
visible depletion events; they do not correspond to any published biophysical
measurement. Simulation results confirm that nanosensors involved in more
detection and gossip events deplete their energy faster, and that the BmIDSet
mechanism faithfully tracks the spread of biomarker information through the
nanosensor network.
\end{abstract}

\begin{IEEEkeywords}
Molecular communication, nanosensor networks, energy modelling, AcCoRD,
reaction-diffusion simulation, Internet of BioNanoThings, biomarker detection,
gossip protocol.
\end{IEEEkeywords}

\IEEEpeerreviewmaketitle

% ============================================================
\section{Introduction}
\label{sec:intro}
% ============================================================

\IEEEPARstart{T}{he} Internet of BioNanoThings (IoBNT) envisions a body-area
network of sub-micron devices embedded in biological fluids that cooperate to
sense, compute, and communicate health-relevant information \cite{akyildiz2015internet}.
Within this vision, nanosensors are the fundamental building blocks: tiny nodes
that drift or adhere in blood or tissue, detect specific molecules (biomarkers)
circulating in the fluid, and relay detection events to neighbouring nodes or
to an external gateway.

A central challenge in realising such networks is energy. Unlike conventional
wireless nodes that can harvest radio-frequency power or carry a battery,
intra-body nanosensors must subsist on energy stored in a nanoscale reservoir
(e.g., a piezoelectric nanogenerator or a biochemical fuel cell
\cite{jornet2012nano}). The reservoir replenishes slowly and is consumed by
every computational and communication operation. If a nanosensor runs out of
energy before it has completed its sensing task, information is lost and the
system silently degrades. Modelling this phenomenon accurately requires a
simulation environment that couples the physical diffusion of molecules, the
stochastic chemistry of biomarker-nanosensor binding, the molecular signalling
between nanosensors (gossip), and the energy bookkeeping of each individual node.

Molecular communication (MC) simulation has matured considerably over the past
decade. Tools such as AcCoRD \cite{noel2017accord} provide accurate stochastic descriptions of diffusion
and first- or second-order chemical reactions. However,
it does not include a per-node energy model that is coupled to
reaction events. Energy is, at best, an offline post-processing concern and
not an in-loop constraint that affects whether a given node can react.

The main contributions of this paper are:
\begin{enumerate}
  \item An individual energy reservoir attached to every simulated nanosensor,
        governed by a passive background drain and event-based costs for
        detection and gossip reactions.
  \item A \emph{biomarker identity set} (BmIDSet) carried by each molecule,
        enabling the simulation to track which unique biomarkers each nanosensor
        has detected and preventing double-counting.
  \item An extension of the AcCoRD JSON configuration format with energy
        parameters that can be populated from published biophysical energy
        models.
  \item A demonstration on a five-region arteriole scenario with per-nanosensor
        energy trajectories, depletion events, and information-spread dynamics.
\end{enumerate}

The remainder of the paper is structured as follows. Section~\ref{sec:background}
introduces the background concepts. Section~\ref{sec:related} reviews related
work. Section~\ref{sec:methodology} details the implementation.
Section~\ref{sec:results} presents and discusses simulation results.
Section~\ref{sec:conclusions} concludes with limitations and future directions.

% ============================================================
\section{Background}
\label{sec:background}
% ============================================================

\subsection{Diffusion-Based Molecular Communication}

In diffusion-based MC, information is encoded in the type, concentration, or
timing of \emph{information molecules} (IMs) released by a transmitter into a
fluid medium \cite{nakano2013molecular}. IMs propagate by Brownian
motion---a thermally-driven random walk that requires no additional energy from
the emitter---and are detected by a receiver nanosensor when they come within
a \emph{binding radius} $r_b$ of the receiver surface. Detection is typically
modelled as a bimolecular reaction:
\begin{equation}
  \mathrm{BM} + \mathrm{NM} \;\xrightarrow{k}\; \mathrm{NM}_{\mathrm{det}},
  \label{eq:detection}
\end{equation}
where BM denotes a biomarker molecule and NM denotes a molecule representing
the receptor state of the nanosensor. The reaction rate $k$ is set high enough
(approaching infinity) to model perfect absorption upon contact.

After detection, nanosensors may wish to alert their neighbours. This is modelled
as a \emph{gossip} reaction, where an informed nanosensor converts a neighbouring
naive nanosensor to an informed state upon close contact:
\begin{equation}
  \mathrm{NM} + \mathrm{NM}_\mathrm{det}
  \;\xrightarrow{k_g}\;
  \mathrm{NM}_\mathrm{det} + \mathrm{NM}_\mathrm{gossip}.
  \label{eq:gossip}
\end{equation}
This reaction propagates awareness of a detection event through the network
without any electromagnetic infrastructure.

\subsection{AcCoRD Simulator}

AcCoRD \cite{noel2017accord} is an open-source, event-driven, hybrid
micro-/mesoscopic simulator for 3-D reaction-diffusion systems. Its key
abstractions are:

\begin{itemize}
  \item \textbf{Regions}: 3-D volumes (rectangular boxes or spheres) simulated
        either in the fine-grained \emph{microscopic} regime (continuous Brownian
        motion, time step $\Delta t$) or the coarser \emph{mesoscopic} regime
        (well-mixed compartments, Gillespie algorithm). Regions can be nested
        and connected.
  \item \textbf{Actors}: logical entities that either release molecules
        (\emph{active} actors) or observe them (\emph{passive} actors). Multiple
        actors can overlap a region.
  \item \textbf{Chemical reactions}: specified globally or per-region, supporting
        zeroth-, first-, and second-order kinetics. Bimolecular reactions in the
        microscopic regime use a binding radius $r_b$.
\end{itemize}

The simulator reads a JSON configuration file and writes observations and a
machine-readable summary to output text files. All stochastic processes are
seeded with a user-controllable random seed, enabling reproducible experiments.

\subsection{Energy Constraints in Nanomachines}

Kuran \emph{et al.}\ \cite{kuran2010energy} derived a model for
MC emission based on the cellular exocytosis pathway. Emission involves four
stages: (1)~molecular synthesis (dominant cost, proportional to peptide
length); (2)~vesicle formation; (3)~vesicle transport along microtubules;
and (4)~membrane fusion via SNARE proteins. Using 1~ATP $\approx 83$~zJ
\cite{freitas1999nanomedicine}, the total emission energy for $n$ insulin
molecules is:
\begin{equation}
  E_T = n E_S + \left\lceil \frac{n}{c_v} \right\rceil (E_V + E_C + E_E),
  \label{eq:kuran}
\end{equation}
where $E_S \approx 10$~aJ per molecule (synthesis cost), and the vesicle
overhead $E_V + E_C + E_E \approx 52.7$~aJ per vesicle (capacity
$c_v \approx 1540$) is negligible by comparison. At cell-scale distances
(1--32~$\mu$m), the total emission energy per symbol ranges from 0.1 to
500~fJ. This makes MC vastly more energy-efficient than electromagnetic
alternatives at intravascular scales.

Importantly, no published model quantifies the energy cost of
\emph{biomarker detection} (receptor binding, signal transduction) or of
\emph{inter-nanosensor gossip}. This is the gap that the AcCoRD energy
extension begins to address by enabling their simulation with configurable
cost parameters.

% ============================================================
\section{Related Work}
\label{sec:related}
% ============================================================

\subsection{Molecular Communication Simulators}

Several simulators have been proposed for MC research.
AcCoRD \cite{noel2017accord} distinguishes itself by combining: (i)~particle-level
microscopic simulation for accuracy in regions of interest; (ii)~mesoscopic
compartment-based simulation for computational efficiency elsewhere;
(iii)~a flexible actor abstraction for modulation schemes including burst,
concentration shift keying (CSK), and point sources; and (iv)~support for
advective flow in addition to diffusion. It has been used to study
a priori Monte Carlo absorption \cite{noel2017accord}, bimolecular reactions
with binding radius, and hybrid regime transitions.

\subsection{Energy Models for Nanosensor Networks}

The energy model of Kuran \emph{et al.}~\cite{kuran2010energy} remains the
reference for MC emission costs. Subsequent works have studied energy
harvesting through piezoelectric effects \cite{jornet2012nano} and biochemical
fuel cells. Akyildiz \emph{et al.}~\cite{akyildiz2015internet} discuss energy
requirements for IoBNT devices in broad terms, noting that energy harvesting and
ultra-low-power operation are essential. However, no closed-form model for
detection or gossip communication energy has been published.

\subsection{Simulators with Energy Modelling}

To the best of our knowledge, no existing MC simulator integrates per-node
energy bookkeeping as an in-loop constraint. Post-processing approaches
(analytically computing expected energy costs from logged reaction counts)
cannot capture the feedback between energy depletion and system behaviour:
when a node runs out of energy, it stops reacting, which affects both the
information it propagates and the reactions of its neighbours. Only an in-loop
energy model can correctly capture this coupling.

% ============================================================
\section{Methodology}
\label{sec:methodology}
% ============================================================

This section describes the energy-aware extensions added to AcCoRD. All
modifications are confined to the source files \texttt{actor.h/c},
\texttt{accord.c}, \texttt{micro\_molecule.h/c}, \texttt{chem\_rxn.h}, and
\texttt{file\_io.c}. The JSON configuration format is extended with new
optional fields; existing configurations are fully backward-compatible.

\subsection{Energy Reservoir Model}

Each AcCoRD actor can be declared energy-enabled with the configuration flag
\texttt{"Is Energy Enabled?"}: \texttt{true}. Six scalar parameters govern the
reservoir:

\begin{itemize}
  \item $E_\mathrm{init}$: initial energy at the start of each repeat
        (\texttt{"Energy Initial"}).
  \item $E_\mathrm{max}$: capacity cap; harvested energy above this is
        discarded (\texttt{"Energy Max"}).
  \item $\dot{E}_\mathrm{drain}$: passive energy drain per global microscopic
        time step $\Delta t$ (\texttt{"Energy Drain Passive"}).
  \item $\dot{E}_\mathrm{harvest}$: passive energy harvest per step
        (\texttt{"Energy Harvest Passive"}).
  \item $e_\mathrm{det}$: energy cost per detection reaction event.
  \item $e_\mathrm{comm}$: energy cost per gossip/communication reaction event.
\end{itemize}

The discrete-time energy update for a single nanosensor unit $u$ is:
\begin{equation}
  E_u(t + \Delta t) =
    \min\!\bigl(E_u(t) - \dot{E}_\mathrm{drain} + \dot{E}_\mathrm{harvest},\;
               E_\mathrm{max}\bigr),
  \label{eq:passive}
\end{equation}
applied at every global microscopic step, plus instantaneous deductions of
$e_\mathrm{det}$ or $e_\mathrm{comm}$ whenever a reaction attributable to
unit $u$ fires.

\subsection{Per-Unit (Per-Nanosensor) Tracking}

A single AcCoRD actor can represent a group of $N$ physically co-located
nanosensors. In the BURST modulation scheme, $N$ molecules are emitted at
initialisation, each molecule representing one physical nanosensor drifting
through the simulation space.

Each actor maintains parallel arrays of length $N$:
\begin{itemize}
  \item \texttt{energyCurrentPerUnit[u]}: current energy of unit $u$.
  \item \texttt{bEnergyDepletedPerUnit[u]}: flag set when
        $E_u < \varepsilon_\mathrm{min}$, where $\varepsilon_\mathrm{min}$ is
        the cost of the least expensive energy-enabled reaction in the
        simulation (i.e.\ the gossip cost $e_\mathrm{comm}$ in the current
        scenario). This threshold is computed automatically at simulation
        start-up by scanning all energy-enabled reactions; it is therefore
        configuration-driven and requires no manual adjustment.
\end{itemize}

The rationale for this threshold is operational: a nanosensor unit whose
remaining energy is strictly below $\varepsilon_\mathrm{min}$ cannot afford
even the cheapest available reaction. Marking it as depleted at that point
prevents the simulator from futilely attempting and aborting every
subsequent reaction involving that unit. Concretely, if the gossip cost is
$e_\mathrm{comm} = 2\times10^{-5}$~arb., a unit with $E_u < 2\times10^{-5}$
is considered depleted. The aggregate \texttt{energyCurrent} is the sum of
all per-unit levels and is used for display purposes only; all energy decisions
operate on individual units.

When a bimolecular reaction involves a molecule owned by unit $u$ of an
energy-enabled actor, the cost is charged to unit $u$. If a unit is depleted,
its molecule cannot participate in energy-gated reactions and is effectively
inert. If all $N$ units of an actor are depleted, the actor's next-action
time is set to infinity, removing it from the simulation event queue.

\subsection{Molecule Ownership}

To attribute energy costs correctly, every molecule in the microscopic regime
carries two additional fields:
\begin{itemize}
  \item \texttt{ownerActorID} (16-bit integer): index into the global actor
        array of the energy-enabled actor that emitted this molecule; $-1$ if
        not energy-owned (e.g., a biomarker or reaction product with no
        nanosensor lineage).
  \item \texttt{ownerUnitID} (32-bit integer): index within the actor's unit
        array.
\end{itemize}
When an active actor emits $N$ molecules in one burst event, molecules are
assigned unit IDs $0, 1, \ldots, N-1$ via a round-robin scheme:
$\texttt{ownerUnitID} = i \bmod N$ for the $i$-th molecule. Ownership is
preserved through reaction products: the product inherits the
\texttt{ownerActorID} and \texttt{ownerUnitID} of the nanosensor reactant.

\subsection{Biomarker Identity Tracking (BmIDSet)}

A key challenge in nanosensor network simulation is preventing double-counting
of the same biomarker. If nanosensor $A$ detects biomarker $b$, gossips to
nanosensor $B$, and $B$ later encounters the same biomarker $b$ again, $B$
should not record another unique detection.

To address this, we introduce the \texttt{BmIDSet} data type: a compact set
of up to 64 globally unique 64-bit biomarker identifiers (IDs) carried by each
molecule. Biomarker molecules are assigned a unique ID at creation via an
atomic counter (\texttt{assignNewBMID()}). Nanosensor molecules begin with an
empty BmIDSet.

When a detection or gossip reaction fires between molecules with BmIDSets $A$
and $B$, the product molecule receives $A \cup B$ (the set union). Each
nanosensor molecule's BmIDSet thus monotonically accumulates the identities of
all biomarkers encountered, directly or via gossip. The simulator logs this
set at each reaction event, enabling reconstruction of the information-spread
graph.

Two molecules may react only when their BmIDSets have no overlap (a guard
implemented in the microscopic reaction checker), preventing a single biomarker
from triggering multiple detection events at the same nanosensor.

\subsection{Energy-Enabled Chemical Reactions}

The AcCoRD reaction structure is extended with four new fields:
\begin{itemize}
  \item \texttt{bEnergyEnabled}: whether this reaction incurs energy costs.
  \item \texttt{energyCostType}: string, \texttt{"detection"} or
        \texttt{"communication"}.
  \item \texttt{energyCostValue}: scalar cost per event.
  \item \texttt{energyCostUnit}: \texttt{"per\_reaction"} (flat cost) or
        \texttt{"per\_product"} (cost scales with number of products).
\end{itemize}

Before a candidate reaction fires, the function
\texttt{canAndConsumeReactionEnergy()} is called. It:
\begin{enumerate}
  \item identifies all energy-enabled actors whose molecules are reactants;
  \item checks whether each identified actor has a unit with sufficient energy;
  \item if yes, atomically deducts the cost from the selected unit of each
        actor and proceeds;
  \item if no, suppresses the reaction.
\end{enumerate}
This \emph{energy gate} means that a nanosensor without energy cannot detect
or communicate, coupling energy state directly to reaction dynamics. In the
microscopic regime, the unit charged is always the one whose molecule physically
participated in the reaction, identified by the \texttt{ownerUnitID} field
carried by every nanosensor molecule. A round-robin fallback is used only for
mesoscopic reactions, where per-molecule ownership is not tracked.

\subsection{Reaction Event Logging}

For every detection and gossip reaction that fires, the simulator logs, per
nanosensor unit:
\begin{itemize}
  \item simulation time $t$;
  \item the BmIDSet of the unit at that moment;
  \item the reaction direction (\texttt{send}, \texttt{receive}, or
        \texttt{send \& receive});
  \item the partner actor ID and unit ID.
\end{itemize}
Duplicate events (same time, same partner) within one time step are suppressed.
The per-unit log is serialised to JSON in the summary file alongside aggregated
statistics: mean energy, final energy, depletion flag, depletion threshold,
total detection count, and total gossip count.

\subsection{MATLAB Analysis Tool}

A companion MATLAB script, \texttt{plotNanomachineEnergy.m}, reads an AcCoRD
output file, parses the per-unit energy time series, and produces one subplot
per actor showing all nanosensor energy curves over time. Depleted units are
highlighted in red. A horizontal dashed line marks the depletion threshold
$\varepsilon_\mathrm{min}$ on each subplot, making it immediately visible which
non-depleted nanosensors are approaching exhaustion. The script requires no
additional toolboxes.

\subsection{Simulation Scenario}

\begin{table}[!t]
\renewcommand{\arraystretch}{1.3}
\caption{Simulation Parameters. Energy values (starred) are arbitrary
placeholders chosen to exercise the model; they carry no physical meaning.}
\label{tab:params}
\centering
\begin{tabular}{lll}
\hline
\textbf{Parameter} & \textbf{Value} & \textbf{Unit} \\
\hline
Global time step $\Delta t$ & $10^{-5}$ & s \\
Simulation duration $T$ & $0.04$ & s \\
Number of regions & 5 & -- \\
Number of molecule types & 4 & -- \\
Diffusion coeff.\ BM ($D_\mathrm{BM}$) & $2.26 \times 10^{-12}$ & m$^2$/s \\
Diffusion coeff.\ NM ($D_\mathrm{NM}$) & $5.66 \times 10^{-14}$ & m$^2$/s \\
BM advection velocity & $1.5 \times 10^{-3}$ & m/s \\
NM advection velocity & $0.5 \times 10^{-3}$ & m/s \\
Number of biomarker sources & 56 & -- \\
Nanosensor actors & 4 & -- \\
Nanosensors per actor ($N$) & 25 & -- \\
Total nanosensors & 100 & -- \\
Binding radius, detection $r_b$ & $1.025$ & $\mu$m \\
Binding radius, gossip $r_g$ & $2.0$ & $\mu$m \\
Initial/max energy $E_\mathrm{init}^*$ & $3 \times 10^{-4}$ & arb. \\
Passive drain $\dot{E}_\mathrm{drain}^*$ & $5 \times 10^{-9}$ & arb./step \\
Detection cost $e_\mathrm{det}^*$ & $9 \times 10^{-5}$ & arb. \\
Gossip cost $e_\mathrm{comm}^*$ & $2 \times 10^{-5}$ & arb. \\
Depletion threshold $\varepsilon_\mathrm{min}^*$ & $2 \times 10^{-5}$ & arb. \\
\hline
\end{tabular}
\end{table}

The scenario models a branching arteriole composed of five rectangular
microscopic regions. The arteriole splits at one junction and reconverges
at another, creating two parallel flow paths. All regions are simulated in
the microscopic regime with a uniform subvolume base size of 1~$\mu$m.

Biomarkers (molecule type 0) are advected at $1.5 \times 10^{-3}$~m/s and
diffuse with coefficient $D_\mathrm{BM} = 2.26 \times 10^{-12}$~m$^2$/s.
A total of 56 biomarker molecules are released at $t=0$ from point sources
distributed along the inlet face. Nanosensors (molecule type 1) are placed
as thin sheets on the vascular wall in four clusters (actors 56--59), each
emitting 25 individual nanosensor molecules at $t=0$.

Four molecule types model the nanosensor state machine:
\begin{enumerate}
  \item BM ($D = 2.26\times10^{-12}$~m$^2$/s, advected): biomarker.
  \item NM ($D = 5.66\times10^{-14}$~m$^2$/s, advected): naive nanosensor.
  \item NM$_\mathrm{det}$ (same diffusivity): nanosensor that has detected a BM.
  \item NM$_\mathrm{gossip}$ (same diffusivity): nanosensor informed via gossip.
\end{enumerate}
Eight bimolecular reactions implement all state transitions. The three
detection reactions (e.g., BM+NM$\to$NM$_\mathrm{det}$) use binding radius
$r_b = 1.025$~$\mu$m; the five gossip reactions use $r_g = 2.0$~$\mu$m.
Binding rates are set to $k = 10^{9999}$ (effectively instantaneous upon
contact). All simulation parameters are summarised in Table~\ref{tab:params}.

\emph{Note on energy values.} The energy constants
($E_\mathrm{init}$, $\dot{E}_\mathrm{drain}$, $e_\mathrm{det}$,
$e_\mathrm{comm}$) were chosen arbitrarily solely to verify that the model
behaves correctly and produces observable depletion within the simulation
window. They do not correspond to any published biophysical measurement and
should not be interpreted as physically realistic. The depletion threshold
$\varepsilon_\mathrm{min} = e_\mathrm{comm} = 2\times10^{-5}$~arb.\ is derived
automatically from the configuration and will update whenever the reaction
costs are changed. Calibrating these parameters to real nanosensor hardware is
left as future work.

% ============================================================
\section{Results}
\label{sec:results}
% ============================================================

\subsection{Energy Depletion}

Table~\ref{tab:depletion} summarises the depletion statistics at the end of
the 40~ms simulation. Out of 100 nanosensors across four actors, twelve were
depleted. No actor lost all its units, meaning that each cluster retained some
detection and gossip capability throughout the simulation.

Actor 56 had the highest depletion rate (6 out of 25 units, 24\%). This actor
is placed near the inlet of the arteriole, where biomarker concentration is
highest and reaction rates are therefore greatest. Furthermore, once a biomarker
reacts with a nanosensor it is consumed, so fewer biomarkers reach downstream
clusters. Actor 57, placed downstream, had a depletion rate of 4 out of 25
(16\%). Actors 58 and 59, located furthest from the inlet along the two
parallel flow paths, each lost only 1 nanosensor (4\%). The position of each
actor cluster relative to the biomarker flow thus has a direct impact on energy
lifetime.

\begin{table}[!t]
\renewcommand{\arraystretch}{1.3}
\caption{Nanosensor Energy Depletion Summary (single run, seed = 5)}
\label{tab:depletion}
\centering
\begin{tabular}{cccc}
\hline
\textbf{Actor ID} & \textbf{NMs} & \textbf{Depleted} & \textbf{Depl.\ rate} \\
\hline
56 & 25 & 6 & 24\% \\
57 & 25 & 4 & 16\% \\
58 & 25 & 1 &  4\% \\
59 & 25 & 1 &  4\% \\
\hline
\end{tabular}
\end{table}

\subsection{Activity--Energy Correlation}

Table~\ref{tab:units} lists four representative nanosensors from actor 56.
The correlation between reaction activity and energy depletion is clear.
NM~2, which participated in no reactions, retained a final energy of
$2.80\times10^{-4}$~arb., close to its initial value (the small deficit is
due to passive drain). NM~22 (1 detection, 1 gossip) was not depleted.
NM~11 (2 detections, 5 gossips) and NM~7 (3 detections, 1 gossip) were both
fully depleted, illustrating that depletion can result from either many
low-cost gossip events or fewer high-cost detection events. Because the energy
values are arbitrary placeholders, no attempt is made to reconcile the exact
numerical differences with a physical model; what matters is that the
qualitative ordering is correct: more active nanosensors lose energy faster,
and the simulation enforces this in-loop so that depleted nanosensors can no
longer react.

\begin{table}[!t]
\renewcommand{\arraystretch}{1.3}
\caption{Representative Nanosensors of Actor 56 (seed = 5)}
\label{tab:units}
\centering
\begin{tabular}{ccccc}
\hline
\textbf{NM} & \textbf{Detections} & \textbf{Gossips} & $E_\mathrm{final}$ (arb.) & \textbf{Depleted?} \\
\hline
NM 2  & 0 & 0 & $2.80\times10^{-4}$ & No  \\
NM 22 & 1 & 1 & $1.70\times10^{-4}$ & No  \\
NM 11 & 2 & 5 & $0$ & Yes \\
NM 7  & 3 & 1 & $0$ & Yes \\
\hline
\end{tabular}
\end{table}

\subsection{BmIDSet Propagation}

Table~\ref{tab:bmids} traces the BmIDSet evolution of NM~6 (actor~56):
starting from an empty set, it first received BM\#26 via gossip, then
directly detected BM\#34, \#22 and \#33 across two events, propagating its
accumulated knowledge to neighbours after each step before becoming depleted.
The absence of any repeated biomarker ID in any unit's log confirms the
deduplication property of the BmIDSet union mechanism.

\begin{table}[!t]
\renewcommand{\arraystretch}{1.3}
\caption{BmIDSet Evolution of NM 6 (Actor 56, seed = 5)}
\label{tab:bmids}
\centering
\begin{tabular}{clll}
\hline
\textbf{Time (ms)} & \textbf{Event} & \textbf{Partner} & \textbf{BmIDSet} \\
\hline
13.71 & gossip $\leftarrow$ & NM 14 & \{26\} \\
14.16 & detect             & BM\#34, \#22 & \{26, 34, 22\} \\
14.17 & gossip $\to$       & NM 18 & \{26, 34, 22\} \\
14.18 & gossip $\to$       & NM 14 & \{26, 34, 22\} \\
14.41 & detect             & BM\#33 & \{26, 34, 22, 33\} \\
14.42 & gossip $\to$       & NM 14 & \{26, 34, 22, 33\} \\
14.43 & gossip $\to$       & NM 18 & \{26, 34, 22, 33\} \\
\hline
\end{tabular}
\end{table}

\subsection{Discussion}

\textbf{Energy parameter scale.}
All energy values used in this experiment are entirely arbitrary placeholders.
They were chosen only to make depletion events observable within the simulation
window and to verify that the energy-gating logic operates correctly. They
carry no physical meaning and should not be compared against the fJ-scale
emission costs of the Kuran model \cite{kuran2010energy}. No biophysical
measurement of nanosensor detection energy or gossip communication energy
currently exists in the literature. The AcCoRD energy extension is designed
to accept any value.

\textbf{Binding radius and molecular geometry.}
The two binding radii used in the scenario reflect the physical sizes of the
molecules involved in each reaction type. In AcCoRD's microscopic regime, the
binding radius $r_b$ represents the centre-to-centre distance below which two
molecules are considered to have reacted; for a bimolecular reaction between
species of radii $r_1$ and $r_2$, the natural choice is $r_b \approx r_1 + r_2$.
For the detection reactions (BM\,+\,NM), the biomarker is small compared to the
nanosensor, yielding $r_b = 1.025~\mu$m. For the gossip reactions
(NM\,+\,NM$_\mathrm{det}$), both reactants are nanosensor-sized molecules, so
their combined radii give a larger binding distance $r_g = 2.0~\mu$m. This
choice ensures that the reaction geometry is consistent with the relative sizes
of the interacting species rather than being an arbitrary tuning parameter.

\textbf{Single gossip per biomarker identity.}
The BmIDSet overlap guard prevents two nanosensor molecules from triggering a
gossip reaction if they already share at least one common biomarker ID. This
ensures that once NM~$A$ has transmitted knowledge of biomarker $b$ to NM~$B$,
a subsequent physical encounter between the same pair does not produce a
redundant and energy-consuming re-transmission of information that
NM~$B$ already holds.

This is a deliberate modelling choice, but it may not reflect the physical
reality of molecular communication. In MC, nanosensor molecules diffuse
freely and can repeatedly come within binding radius of one another; without
an explicit molecular memory (e.g.\ a conformational change or covalent
modification that marks the molecule as ``already informed''), there is no
biophysical mechanism to suppress re-reaction. A more conservative model
would allow every encounter to trigger a gossip reaction regardless of
prior history, which would increase energy consumption without adding new
information to the network. Whether the BmIDSet guard is a reasonable
abstraction or an over-idealisation depends on the nanosensor architecture
assumed; this remains an open modelling question for future work.

\textbf{Depletion and information loss.}
NM~6 (2 detections, 5 gossips) depleted fully and ceased to participate in the
reaction system. NM~7 reached the same outcome through a different path: three
detections and only one gossip event sufficed to exhaust its reservoir, because
each detection costs $9\times10^{-5}$~arb.\ compared to $2\times10^{-5}$~arb.\
for gossip. Any subsequent biomarkers reaching these nanosensors' locations were
not detected, and they no longer propagated gossip to their neighbours. This
information-loss effect cannot be reproduced by any post-processing energy model
and demonstrates the necessity of an in-loop energy gate.

% ============================================================
\section{Conclusions}
\label{sec:conclusions}
% ============================================================

We have presented an energy-aware extension of AcCoRD that integrates
individual nanosensor energy reservoirs directly into the simulation event loop.
The extension introduces passive energy drain, event-based costs for biomarker
detection and inter-nanosensor gossip, a biomarker identity set (BmIDSet) for
information tracking and deduplication, and a full per-unit logging
infrastructure. The JSON configuration interface is backward-compatible with
all existing AcCoRD simulations, and a MATLAB post-processing script enables
straightforward visualisation of energy trajectories.

Simulation of a branching arteriole scenario with 100 nanosensors and 56
biomarkers confirms the expected behaviour: nanosensors with higher detection
and gossip activity deplete faster, the BmIDSet mechanism faithfully tracks
information propagation, and energy depletion silently removes exhausted nodes
from the system dynamics in a physically meaningful way.

\subsection{Limitations}

\begin{itemize}
  \item The simulation was run for a single random seed. Statistical analysis
        over many realisations is needed to characterise average depletion rates
        and information coverage.
  \item All energy parameter values are arbitrary placeholders used to test
        the model; no biophysical measurement of detection or gossip costs
        currently exists, and the simulation results cannot be interpreted as
        physically realistic energy budgets.
  \item The BmIDSet capacity is fixed at 64 unique biomarker IDs per molecule.
        Simulations with more than 64 distinct biomarker species require an
        enlarged or dynamically-allocated structure.
  \item Molecule ownership and BmIDSet tracking are limited to the microscopic
        regime; mesoscopic compartments do not yet carry these attributes.
\end{itemize}

\subsection{Future Work}

Several directions extend this work naturally. First, a parametric study
varying $e_\mathrm{det}$, $e_\mathrm{comm}$, and $\dot{E}_\mathrm{drain}$
across physically motivated ranges---informed by the Kuran emission model
\cite{kuran2010energy}---would establish which parameters most strongly affect
network lifetime and detection coverage. Second, event-triggered energy
harvesting (e.g., from glucose oxidation when the nanosensor encounters a
sugar molecule) could complement the existing constant-rate harvest parameter.
Third, the BmIDSet mechanism could implement distributed detection algorithms
such as $k$-detection consensus, where a nanosensor alerts only after
accumulating evidence of $k$ distinct biomarker encounters. Finally, extending
ownership and BmIDSet tracking to the mesoscopic regime would enable
energy-lifetime analysis across multi-scale hybrid simulations.

\appendices

\section{JSON Configuration Extensions}
\label{app:json}

The following fields are added to each actor entry in the
\texttt{"Actor Specification"} array. All fields are optional; omitting them
disables the energy system for that actor.

\begin{itemize}
  \item \texttt{"Is Energy Enabled?"}: boolean (default \texttt{false}).
  \item \texttt{"Energy Initial"}: scalar (default 0).
  \item \texttt{"Energy Max"}: scalar (default 0).
  \item \texttt{"Energy Drain Passive"}: scalar per global time step (default 0).
  \item \texttt{"Energy Harvest Passive"}: scalar per global time step (default 0).
\end{itemize}

The following fields are added to each entry in the
\texttt{"Chemical Reaction Specification"} array:

\begin{itemize}
  \item \texttt{"Is Energy Enabled?"}: boolean (default \texttt{false}).
  \item \texttt{"Energy Cost Type"}: string, \texttt{"detection"} or
        \texttt{"communication"}.
  \item \texttt{"Energy Cost Value"}: scalar (default 0).
  \item \texttt{"Energy Cost Unit"}: \texttt{"per\_reaction"} or
        \texttt{"per\_product"} (default \texttt{"per\_reaction"}).
\end{itemize}

The depletion threshold $\varepsilon_\mathrm{min}$ is not a configuration
parameter: it is derived automatically at simulation start-up as the minimum
\texttt{"Energy Cost Value"} across all energy-enabled reactions, so it
always stays consistent with the reaction costs without any manual adjustment.

\section*{Acknowledgment}

The author thanks the developers of AcCoRD (Adam Noel \emph{et al.}) for
making the simulator open-source, and the I2R laboratory for providing
computational resources and research support.

\ifCLASSOPTIONcaptionsoff
  \newpage
\fi

\begin{thebibliography}{7}

\bibitem{akyildiz2015internet}
I.~F. Akyildiz, M.~Pierobon, S.~Balasubramaniam, and Y.~Koucheryavy,
``The internet of bio-nano things,''
\emph{IEEE Communications Magazine}, vol.~53, no.~3, pp.~32--40, Mar.~2015.

\bibitem{noel2017accord}
A.~Noel, K.~C. Cheung, and R.~Schober,
``Improving the accuracy of diffusion-based molecular communication channels
using ${A}$~priori Monte Carlo absorption,''
in \emph{Proc. IEEE International Conference on Communications (ICC)}, Paris,
France, May 2017.

\bibitem{kuran2010energy}
M.~\c{S}.~Kuran, H.~B. Yilmaz, T.~Tugcu, and I.~F. Akyildiz,
``Signalling and communication energy efficiency for two-way communication in
molecular nanonetworks,''
in \emph{Proc. IEEE INFOCOM}, San Diego, CA, Mar.~2010, pp.~1--9.

\bibitem{jornet2012nano}
J.~M. Jornet and I.~F. Akyildiz,
``Joint energy harvesting and communication analysis for perpetual wireless
nanosensor networks in the terahertz band,''
\emph{IEEE Transactions on Nanotechnology}, vol.~11, no.~3, pp.~570--580,
May~2012.

\bibitem{freitas1999nanomedicine}
R.~A. Freitas,
\emph{Nanomedicine, Volume~I: Basic Capabilities}.
Georgetown, TX: Landes Bioscience, 1999.

\bibitem{nakano2013molecular}
T.~Nakano, A.~W. Eckford, and T.~Haraguchi,
\emph{Molecular Communication}.
Cambridge, UK: Cambridge University Press, 2013.

\bibitem{tiwari2020molsim}
R.~Tiwari and A.~Bhatt,
``MolSim: A MATLAB-based simulation framework for molecular communication,''
in \emph{Proc. IEEE GLOBECOM Workshops}, Taipei, Taiwan, Dec.~2020, pp.~1--6.

\end{thebibliography}



\end{document}
